\section{Ray-Objekt Schnitte}

In diesem Kapitel werden Schnitttests mit unterschiedlichen Objekten vorgestellt. Die allgemeine
Vorgehensweise besteht darin, für Objekte eine Oberfläche in implizierter Form zu definieren.

\begin{equation}
    f(x, y, z) = f(\mathbf{p}) = 0
    \label{lb_surface}
\end{equation}

In dieser Form lässt sich prüfen ob ein Punkt auf der Oberfläche liegt, indem er in die Funktion eingesetzt wird
und diese "0" ergeben muss. Ist das Ergebnis von "0" verschieden so befindet sich der Punkt oberhalb oder unterhalb der Oberfläche.

Durch die parametrische Darstellung eines Rays $\mathbf{r}$ (Strahls) lassen sich alle Punkte entlang des Strahls
durch die geeignete Wahl das Parameter $\mathit{t}$ berechnen. 

Sei $\mathbf{o}$ ein Stützpunkt in homogenen Koordinaten auf dem Strahl und $\mathbf{d}$ ein Richtungsvektor
in homogenen Koordinaten, dann ist $\mathbf{r}$ definiert durch:
\begin{equation}
    \mathbf{r} = \mathbf{o} + \mathit{t} \cdot \mathbf{d}
    \label{eq_ray}
\end{equation}


Man setzt nun Gleichung~\ref{eq_ray} in Gleichung~\ref{lb_surface} ein und löst nach $\mathit{t}$ auf.

\begin{equation}
    f(\mathbf{o} + \mathit{t} \cdot \mathbf{d}) = 0
\end{equation}

Falls der Ray die Oberfläche schneidet lassen sich konkrete Werte für $\mathit{t}$ berechnen. Dabei
ist man am kleinsten Wert für $\mathit{t}$ interessiert, da man den Schnittpunkt mit der am nächsten zum
Betrachter zugewandten Oberflächenseite berechnen will. Der exakte
Schnittpunkt mit der Oberfläche ergibt sich dann durch einsetzen des für $\mathit{t}$ ermittelten Wertes
in Gleichung~\ref{eq_ray}.


\subsection{Ray-Ebenen Schnitt}

Die implizite Form einer Ebene lautet:
\begin{equation}
    f(\mathbf{p}) = (\mathbf{p} - \mathbf{x}) \cdot \mathbf{n} = 0
\end{equation}

Wobei $\mathbf{x}$ ein Punkt ist der in der Ebene liegt und $\mathbf{n}$ ist ein Normalenvektor, der senkrecht
auf der Ebene steht. Die Gleichung lässt sich nun wie folgt umformen:


\begin{equation*}
    \mathbf{p} \cdot \mathbf{n} = \mathbf{x} \cdot \mathbf{n}
\end{equation*}

Einsetzen von Gleichung~\ref{eq_ray} liefert:

\begin{align*}
    \mathbf{o} \cdot \mathbf{n} + \mathit{t} \cdot \mathbf{d} \cdot \mathbf{n} &= \mathbf{x} \cdot \mathbf{n}\\
    \mathit{t} \cdot \mathbf{d} \cdot \mathbf{n} &= \mathbf{x} \cdot \mathbf{n} - \mathbf{o} \cdot \mathbf{n}\\
    \mathit{t} &= \dfrac{(\mathbf{x}  - \mathbf{o}) \cdot \mathbf{n}}{\mathbf{d} \cdot \mathbf{n}}
\end{align*}

Damit besitzt $\mathit{t}$ eine eindeutige Lösung. Jedoch sollte der Fall, wenn $\mathbf{d}$ und $\mathbf{n}$
senkrecht aufeinander stehen in der Implementierung gesondert behandelt werden, da dies zu einer Division durch
"0" führt.


\subsection{Ray-Rechteck Schnitt}

Um ein Rechteck darzustellen reicht es aus einen Eckpunkt $\mathbf{c}$ (hier linke untere Ecke) festzulegen,
sowie zwei aufeinander senkrecht stehende Vektoren $\mathbf{w}$ und $\mathbf{h}$, die die Breite bzw. Höhe
des Rechtecks festlegen. Durch das Kreuzprodukt $\mathbf{w} \times \mathbf{h} = \mathbf{n}$  wird die Normale 
$\mathbf{n}$ berechnet. Diese wird später für die Beleuchtung gebraucht, spielt aber auch ein tragende Rolle
beim Ray-Rechteck Schnitttest.

%%%%%%%%%%%%%%%%%%%%%%%%%%%%%%%%%%%%Abbildung%%%%%%%%%%%%%%%%%%%%%%%%%%%%%%%%%%%%%%%%%%%%%%%%%%%%%%%%%%%%
\begin{figure}[h]
	\centering
  \includegraphics[width=\textwidth]{images/Ray-Rectangle.pdf}
	\caption{Ein Ray schneidet das Rechteck im Punkt $\mathbf{p}$}
	\label{fig_Ray-Rectangle}
\end{figure}
%%%%%%%%%%%%%%%%%%%%%%%%%%%%%%%%%%%%%%%%%%%%%%%%%%%%%%%%%%%%%%%%%%%%%%%%%%%%%%%%%%%%%%%%%%%%%%%%%%%%%%%%%

Im ersten Schritt muss nämlich der Schnittpunkt $\mathbf{p}$ (siehe Abbildung~\ref{fig_Ray-Rectangle}) 
berechnet werden. Dies geschieht durch einen Ray-Ebenen Schnitttest, wobei die Ebene gegeben ist durch
den Stützpunkt $\mathbf{c}$ und der Normalen $\mathbf{n}$. Anschließend wird der Vektor $\mathbf{d} = 
\mathbf{p} - \mathbf{c}$ auf die Vektoren $\mathbf{w}$ und $\mathbf{h}$ projiziert, wobei 
$\mathbf{\|w\|}$ und $\mathbf{\|h\|}$ die Einheitsvektoren von $\mathbf{w}$ und $\mathbf{h}$ sind:

\begin{align*}
    \mathbf{dw}& = \mathbf{d} \cdot \mathbf{\|w\|}\\
    \mathbf{dh}& = \mathbf{d} \cdot \mathbf{\|h\|}\\
\end{align*}

Nun müssen vier Fälle geprüft werden. Werden alle dieser Fälle wahr ist, so hat der Ray das Rechteck geschnitten
und der Schnittpunkt kann mit Hilfe von $\mathit{t}$ berechnet werden. Wenn nur einer der Test fehlschlägt, dann
liegt kein Schnitt vor. Die Tests sind:

\begin{enumerate}
\centering
\item $\mathbf{|dw|} \geq 0$
\item $\mathbf{|dh|} \geq 0$
\item $\mathbf{|dw|} \leq \mathbf{|w|}$
\item $\mathbf{|dh|} \leq \mathbf{|h|}$
\end{enumerate}



\subsection{Ray-Kugel Schnitt}

Alle Punkte, die auf der Oberfläche einer Kugel liegen erfüllen die implizite Form:

\begin{equation}
    (\mathbf{p} - \mathbf{c})^{2} - \mathit{r}^2 = 0
    \label{eq_sphere}
\end{equation}

Wobei der Punkt $\mathbf{c}$ den Mittelpunkt der Kugel und $\mathit{r}$ den Radius der Kugel
angibt. $\mathbf{p}$ ist ein Punkt der beliebig gewählt werden kann. Jedoch besitzt die
Gleichung~\ref{eq_sphere} nur dann eine Lösung wenn $\mathbf{p}$ auf der Kugeloberfläche liegt.
Ziel ist es also für $\mathbf{p}$ die parametrische Ray Darstellung aus Gleichung~\ref{eq_ray}
einzusetzen und nach dem Parameter $\mathit{t}$ aufzulösen. Nur wenn der Ray die Kugel schneidet
oder tangiert gibt es eine Lösung für $\mathit{t}$.

Bevor Gleichung~\ref{eq_ray} für $\mathbf{p}$ eingesetzt wird erfolgt noch ein Umformungsschritt:
\begin{equation*}
    (\mathbf{p}^2 - 2\mathbf{pc} + \mathbf{c}^2) - \mathit{r}^2 = 0
\end{equation*}

Einsetzen von Gleichung~\ref{eq_ray} liefert:

\begin{align*}
    &(\mathbf{o} + \mathit{t} \cdot \mathbf{d})^2 - 2 \cdot (\mathbf{o} + \mathit{t} \cdot 
      \mathbf{d}) \cdot \mathbf{c} + \mathbf{c}^2 - \mathit{r}^2 = 0\\
    &(\mathbf{o}^2 + 2\mathbf{od}\mathit{t} + (\mathit{t}\mathbf{d})^2) - 
     (2\mathbf{oc} + 2\mathbf{cd}\mathit{t}) + \mathbf{c}^2 - \mathit{r}^2 = 0\\
    &\mathbf{d}^2\mathit{t}^2 - 2\mathbf{cd}\mathit{t} + 2\mathbf{do}\mathit{t} + \mathbf{o}^2 -
     2\mathbf{oc} + \mathbf{c}^2 - \mathit{r}^2 = 0\\
    &\mathbf{d}^2\mathit{t}^2 + 2\mathit{t}\mathbf{d} \cdot (\mathbf{o} - \mathbf{c}) +
     (\mathbf{o} - \mathbf{c})^2 - \mathit{r}^2 = 0
\end{align*}

In der letzten Zeile liegt die Gleichung in Form einer quadratischen Gleichung vor, die bekanntlich
mit Hilfe der Mitternachtsformel gelöst werden kann.

\begin{equation*}
   t = \frac{-b \pm \sqrt{b^2 - 4ac}}{2a} 
\end{equation*}

Die Parameter $\mathit{a, b, c}$ sind dabei wie folgt zu belegen:

\begin{align*}
    a &= \mathbf{d}^2\\
    b &= 2\mathbf{d} \cdot (\mathbf{o} - \mathbf{c})\\
    c &= (\mathbf{o} - \mathbf{c})^2 - \mathit{r}^2
\end{align*}

Für die Implementierung ist es sinnvoll zuerst die Diskriminante zu berechnen und
auf negative Ergebnisse zu prüfen. So kann frühzeitig festgestellt werden, dass
ein Strahl die Kugel nicht schneidet.

\subsection{Ray-Dreieck Schnitt}

Eines der wichtigsten geometrischen Objekte in der Computergrafik ist das Dreieck.
"`Kurvige"' Oberflächen lassen sich zwar gut durch "`NURBS"' darstellen,
jedoch lassen sich, damit Schnitttests nicht in Echtzeit durchführen. Aus diesem Grund
werden die meisten Oberflächen trianguliert als "`triangle meshes"' gespeichert.\\

Für den Ray-Dreieck Schnittest werden Baryzentrische Koordinaten verwendet, diese
beschreiben die relative Lage eines Punktes $\mathbf{p}(\mathit{\alpha, \beta})$ anhand 
der Parameter $\mathit{\alpha, \beta}$ und drei Punkten $\mathbf{p_1, p_2, p_3}$ die das
Dreieck aufspannen.

$\mathit{\alpha, \beta}$ kann nun durch: 

\begin{equation}
    \mathbf{p}(\mathit{\alpha, \beta}) = (1 - \alpha - \beta)\mathbf{p_1} + \alpha\mathbf{p_2} + \beta\mathbf{p_3}
    \label{eq_tri}
\end{equation}

Der Punkt $\mathbf{p}(\mathit{\alpha, \beta})$ liegt innerhalb des Dreiecks
wenn folgende Bedingungen erfüllt sind:


\begin{enumerate}
\centering
\item $\mathit{\alpha} \geq 0$
\item $\mathit{\beta} \geq 0$
\item $\mathit{\alpha} + \mathit{\beta} \leq 1$
\end{enumerate}

Gesucht sind nun die Werte für $\mathit{t}$, $\mathit{\alpha}$ und $\mathit{\beta}$.
Dazu setzt man die parametrische Ray Darstellung~\ref{eq_ray} in Gleichung~\ref{eq_tri} ein:

\begin{align*}
    \mathbf{o} + \mathit{t}\mathbf{d} &= (1 - \alpha - \beta)\mathbf{p_1} + \mathit{\alpha}\mathbf{p_2}
                                          + \mathit{\beta}\mathbf{p_3}\\
    \mathbf{o} + \mathit{t}\mathbf{d} &= \mathbf{p_1} - \mathbf{p_1}\mathit{\alpha} - 
    \mathbf{p_1}\mathit{\beta} + \mathbf{p_2}\mathit{\alpha} + \mathbf{p_3}\mathit{\beta}\\
    \mathbf{o} + \mathit{t}\mathbf{d} &= \mathbf{p_1} + \mathit{\alpha}(\mathbf{p_2 - p_1})
    + \mathbf{\beta}(\mathbf{p_3 - p_1})\\
    \mathbf{o - p_1} &= \mathit{\alpha}(\mathbf{p_2 - p_1}) + \mathbf{\beta}(\mathbf{p_3 - p_1})
    - \mathit{t}\mathbf{d}
\end{align*}


Nun folgenden Substitutionen $\mathbf{s = o - p_1}$, $\mathbf{e_1 = p_2 - p_1}$ und 
$\mathbf{e_2 = p_3 - p_1}$ erhält man das Gleichungssystem

\begin{equation*}
    \mathbf{s} = \begin{pmatrix} \mathbf{-d} && \mathbf{e_1} && \mathbf{e_2} \end{pmatrix} \cdot
    \begin{pmatrix} \mathit{t} \\ \mathit{\alpha} \\ \mathit{\beta} \end{pmatrix}
\end{equation*}

welches man mit Hilfe der Cramerschen Regel lösen kann.
Als Lösung ergibt sich folgender Ausdruck:

\begin{equation}
    \begin{pmatrix} \mathit{t} \\ \mathit{\alpha} \\ \mathit{\beta} \end{pmatrix} =
    \frac{1}{(\mathbf{d \times e_2}) \cdot \mathbf{e_1}} \cdot
    \begin{pmatrix}
        (\mathbf{s \times e_1}) \cdot \mathbf{e_2}\\
        (\mathbf{d \times e_2}) \cdot \mathbf{s}\\
        (\mathbf{s \times e_1}) \cdot \mathbf{d}
    \end{pmatrix}
    \label{eq_tri_result}
\end{equation}